\chapter{Odtworzenie środowiska i eksperymentów}
\label{app:reproduction}
\section{Źródła i wersje}
Repozytorium projektu: \href{\ThesisRepository}{\texttt{github.com/raaddi/edgeguard-iot}}.
\writingnote{Przed oddaniem pracy wskaż konkretny tag lub commit, wersje zależności i instrukcję odtworzenia eksperymentów.}
\section{Budowanie dokumentu}
Aktualne polecenia kompilacji znajdują się w pliku \texttt{thesis/README.md}. Praca offline wymaga wcześniejszego przygotowania kompilatora i używanych pakietów.
\section{Manifest przebiegu}
Poniższy plik pokazuje strukturę planowanych metadanych. Wartości \texttt{null} oznaczają brak wykonanych pomiarów.
\par\noindent\begin{minipage}{\linewidth}
\IfFileExists{../experiments/templates/run-manifest.example.json}{%
  \lstinputlisting[caption={Szablon manifestu, nie wynik eksperymentu.},label={lst:manifest}]{../experiments/templates/run-manifest.example.json}
}{%
  \lstinputlisting[caption={Szablon manifestu, nie wynik eksperymentu.},label={lst:manifest}]{assets/run-manifest.example.json}
}
\end{minipage}\par
\section{Dostęp do danych}
\writingnote{Podaj źródła, licencje i sumy kontrolne danych oraz sposób ich przygotowania. Dużych zbiorów i plików modeli nie umieszczaj bezpośrednio w Git.}
